% TOP_PROBLEM: {"id": 3360, "title": "Giuga's Conjecture on Primality", "category_id": 1, "category": {"id": 1, "name": "number_theory", "display_name": "Number Theory", "description": "Properties of integers, prime numbers, Diophantine equations.", "slug": "number-theory", "order_index": 1, "created_at": "2026-07-31T15:26:25.670Z"}, "status": "open", "problem_number": "OPG-37411", "set_id": 12}
% FIRST_POSTED: 2026-10-01
% ENTRY_KIND: statement_only
% UnsolvedMath ID 3360; source catalogue code OPG-37411
% !TeX program = lualatex
% Definition and sources only; this file is not an LLM solution attempt.
\documentclass[11pt,a4paper]{article}
\usepackage[margin=25mm]{geometry}
\usepackage{fontspec}
\setmainfont{lmroman10-regular.otf}[BoldFont=lmroman10-bold.otf,
 ItalicFont=lmroman10-italic.otf,BoldItalicFont=lmroman10-bolditalic.otf]
\usepackage{amsmath,amssymb,mathtools}
\usepackage{unicode-math}
\setmathfont{latinmodern-math.otf}
\AtBeginDocument{\let\setminus\smallsetminus}
\usepackage{xurl,hyperref}
\hypersetup{colorlinks=true,urlcolor=blue}
\setcounter{secnumdepth}{0}
\setlength{\parindent}{0pt}
\setlength{\parskip}{5pt}
\setlength{\emergencystretch}{3em}
\begin{document}
\section*{Giuga's Conjecture on Primality}\label{3360}
{\small UnsolvedMath ID 3360 · OPG-37411 · Number Theory · OpenGarden}
\subsection{Definitions and mathematical statement}
For an integer $n\ge2$, put $S(n)=\sum_{a=1}^{n-1}a^{n-1}$.
Congruence $S(n)\equiv-1\pmod n$ means that $n$ divides the integer
$S(n)+1$. Giuga's conjecture states
\[
 S(n)\equiv-1\pmod n\quad\Longleftrightarrow\quad n\text{ is prime}.
\]
The restriction $n\ge2$ removes the degenerate modulus one. For a
prime $n$, each nonzero residue satisfies $a^{n-1}\equiv1\pmod n$,
so the forward primality implication is the substantive part:
no composite integer should satisfy the displayed sum congruence.

The exponent is tied to the tested integer, and the sum runs over
every integer from $1$ to $n-1$, including residues not coprime to $n$.
It is therefore not a condition on just the unit group, nor a demand
that each summand separately be congruent to one. The asserted
criterion is exact, without exceptional composites beyond a cutoff.
One can compute the sum using modular exponentiation, but a finite
search by itself proves only a bounded range. A disproof requires
a composite $n$, together with enough factorization and modular
calculation to certify both its compositeness and the sum congruence.
\subsection{Short English statement}
Can a composite integer pass the primality test obtained by summing the (n minus one)st powers of every nonzero residue modulo n?
\subsection{Sources}
\begin{itemize}
\item[S1] ULAM.AI and the UnsolvedMath Contributors, supplied problems.json, record 3360 (OPG-37411), OpenGarden collection. Catalogue identity and source transcription; CC BY 4.0.
\url{https://huggingface.co/datasets/ulamai/UnsolvedMath}.
\item[S2] Open Problem Garden, Giuga's Conjecture on Primality. Original problem and discussion; the mathematical formulation below is expanded from the supplied source record.
\url{http://www.openproblemgarden.org/op/giugas_conjecture_on_primality}.
\end{itemize}
\end{document}
